\section{Prior-Driven Classical Routing Framework}
\label{sec:method}

In this section, we present the mathematical formulation of our Prior-Driven Classical Routing Framework. The central pillar of our framework is the **Zero-Initialized Softmax prior** combined with $L_2$ weight decay, which serves as a powerful implicit regularizer. To theoretically understand and guide this stabilization under heterogeneous batch streams, we introduce a group-level limit constraint, \textbf{Task-Variance Regularization ($\mathcal{L}_{VR}$)}, designed to suppress intra-task routing variance and prevent Heterogeneity Collapse. Crucially, as we show in our empirical audit, our zero-initialized Softmax architectural prior naturally and inherently satisfies this group-level variance limit, making explicit penalty tuning empirically redundant in practice, yet providing high diagnostic value. 

First, we define the mechanics of dynamic model merging in the parameter space. Next, we detail our low-dimensional unit-state feature projection and layer-wise routing projections. Finally, we introduce our multi-objective calibration loss, formalize the group-level \textbf{Task-Variance Regularization ($\mathcal{L}_{VR}$)} constraint, and deconstruct the role of zero-initialized Softmax routing as the foundational implicit prior-layer regularizer.

\subsection{Dynamic Model Merging in Parameter Space}
Let $W_{\text{base}}^{(l)}$ represent the pre-trained base parameters (weights and biases) at layer $l \in \{1, \dots, L\}$. Suppose we have $K$ independent task-specific visual expert models, where the parameters of expert $k \in \{1, \dots, K\}$ at layer $l$ are denoted by $W_k^{(l)}$. Following the task-vector paradigm \cite{ilharco2022editing}, we define the task vector $V_k^{(l)}$ as the parameter difference between the specialized expert and the base model:
\begin{equation}
    V_k^{(l)} = W_k^{(l)} - W_{\text{base}}^{(l)}
\end{equation}

For an incoming batch of inputs $x$ of size $B$, the dynamically merged parameter set $W_{\text{merged}}^{(l)}(x)$ at layer $l$ is assembled by adding a weighted combination of task vectors to the base parameters:
\begin{equation}
    W_{\text{merged}}^{(l)}(x) = W_{\text{base}}^{(l)} + \sum_{k=1}^K \bar{\alpha}_k(l) V_k^{(l)}
\end{equation}
where $\bar{\alpha}_k(l) \in \mathbb{R}$ represents the batch-averaged merging coefficient for task $k$ at layer $l$:
\begin{equation}
    \bar{\alpha}_k(l) = \frac{1}{B} \sum_{b=1}^B \alpha_{k, b}(l)
\end{equation}
Here, $\alpha_{k, b}(l)$ represents the sample-wise merging coefficient computed specifically for input sample $b \in \{1, \dots, B\}$.

\subsection{Low-Dimensional Unit-State Feature Projection}
Let $z(x)_b \in \mathbb{R}^D$ be the globally pooled latent representation extracted from the backbone network's early layers for sample $b$. In deep network sandboxes or large-scale vision models (e.g., CLIP ViT-B/16), the feature dimension $D$ is typically high (e.g., $D=192$ or $D=768$). Direct mapping from this high-dimensional space to routing coefficients on scarce calibration data (e.g., $|D_{\text{cal}}| = 64$) inevitably leads to severe overfitting.

To prevent this, we project $z(x)_b$ into a highly compact $d$-dimensional subspace ($d = K \ll D$) using a static, frozen projection matrix $P \in \mathbb{R}^{D \times d}$. The projection matrix is initialized as a stable, normalized random projection matrix satisfying the Johnson-Lindenstrauss lemma, which provides a data-free and computationally efficient mapping to the $d$-dimensional subspace. To ensure stability against scale variations, we normalize the projected representation onto a unit sphere, obtaining the unit state $\psi(x)_b$:
\begin{equation}
    \psi(x)_b = \frac{z(x)_b P}{\|z(x)_b P\|_2 + \epsilon} \in \mathbb{R}^d
\end{equation}
where $\epsilon = 10^{-8}$ is a small numerical stabilizer.

\subsection{Layer-wise Classical Routing}
For each layer $l \in \{1, \dots, L\}$, we compute the sample-wise merging coefficient $\alpha_{k, b}(l)$ for each task expert $k$ by applying a Softmax activation over the linear projection outputs across the task experts:
\begin{equation}
\begin{aligned}
    \alpha_{k, b}(l) &= \text{Softmax}\left( \langle \psi(x)_b, W_{l, \cdot} \rangle + B_{l, \cdot} \right)_k \\
    &= \frac{\exp\left( \langle \psi(x)_b, W_{l, k} \rangle + B_{l, k} \right)}{\sum_{k'=1}^K \exp\left( \langle \psi(x)_b, W_{l, k'} \rangle + B_{l, k'} \right)}
\end{aligned}
\end{equation}
where $W_{l, \cdot} \in \mathbb{R}^{K \times d}$ and $B_{l, \cdot} \in \mathbb{R}^K$ represent the trainable routing weights and biases for layer $l$ (with $W_{l, k} \in \mathbb{R}^d$ and $B_{l, k} \in \mathbb{R}$ corresponding to task expert $k$). The Softmax normalizes the coefficients across all $K$ tasks to sum to $1$ for each sample.

Unlike quantum-inspired methods (e.g., QWS-Merge \cite{qwsmerge2025}) that pass these projection values through non-monotonic wave-interference activations:
\begin{equation}
    \alpha^{\text{QWS}}_{k, b}(l) = \cos\left( \langle \psi(x)_b, W_{l, k} \rangle + B_{l, k} \right)^2
\end{equation}
our classical projections are linear, avoiding the rugged, non-convex optimization landscapes that trap gradient descent in small-sample calibration splits.

\subsection{Task-Variance Regularization ($\mathcal{L}_{VR}$)}
\label{subsec:vr_loss}

During deployment, inference streams are frequently heterogeneous, meaning they contain batches mixed with samples from different tasks. Standard dynamic routers batch-average sample-wise coefficients to construct a single merged model per batch:
\begin{equation}
    \bar{\alpha}_k(l) = \frac{1}{B} \sum_{b=1}^B \alpha_{k, b}(l)
\end{equation}
When the batch is mixed, this batch-averaging collapses task-specific routing weights toward a uniform, intermediate compromise (e.g., $\bar{\alpha}_k(l) \approx 0.25$ for all $k$). This is \emph{heterogeneity collapse}: the specialized parameters needed for task $A$ are diluted by the requirements of task $B$ within the same batch, severely degrading multi-task accuracy.

To address this and model the optimal group-level behavior of dynamic ensembling, we introduce a mathematical limit constraint, \textbf{Task-Variance Regularization ($\mathcal{L}_{VR}$)}. Although we will show in Section~\ref{sec:experiments} that this explicit loss penalty is empirically redundant once our central zero-initialized Softmax architectural prior is established, formalizing $\mathcal{L}_{VR}$ provides a mathematically precise limit target that characterizes the optimal low-variance group behavior required for robust multi-task ensembling. During the calibration phase, we analyze the router on simulated heterogeneous batches of varying mixtures. Let $S_k \subseteq \{1, \dots, B\}$ represent the indices of samples in a calibration batch that belong to task expert $k$. When $S_k \neq \emptyset$, we compute the intra-task sample variance of predicted routing coefficients at layer $l$:
\begin{equation}
    \sigma^2_{\text{task}}(k, l) = \frac{1}{|S_k|} \sum_{b \in S_k} \left( \alpha_{k, b}(l) - \mu_k(l) \right)^2
\end{equation}
where $\mu_k(l)$ is the mean predicted coefficient for task $k$ across its respective sample group $S_k$:
\begin{equation}
    \mu_k(l) = \frac{1}{|S_k|} \sum_{b \in S_k} \alpha_{k, b}(l)
\end{equation}
In practice, a calibration batch might not contain any samples for a given task, rendering $S_k = \emptyset$. Thus, the task-variance penalty $\mathcal{L}_{VR}$ is averaged only over non-empty (active) task groups, which we denote as $\mathcal{K}_{\text{active}} = \{k \in \{1,\dots,K\} \mid S_k \neq \emptyset\}$:
\begin{equation}
    \mathcal{L}_{VR} = \lambda_{\text{var}} \frac{1}{L \cdot |\mathcal{K}_{\text{active}}|} \sum_{l=1}^L \sum_{k \in \mathcal{K}_{\text{active}}} \sigma^2_{\text{task}}(k, l)
\end{equation}
where $\lambda_{\text{var}} \ge 0$ is a hyperparameter scaling the penalty. Minimizing $\mathcal{L}_{VR}$ forces the router to produce highly consistent, low-variance routing predictions within each task group, suppressing sample-specific noise and outlier skews that lead to heterogeneity collapse. (Note that we explicitly employ the uncorrected population variance formula with factor $1/|S_k|$ rather than Bessel's corrected sample variance with factor $1/(|S_k|-1)$. In the edge case where a calibration batch contains exactly one sample for a task group, i.e., $|S_k| = 1$, Bessel's correction would lead to an undefined division-by-zero error, whereas our formulation naturally evaluates to exactly zero for that group, ensuring stable numerical gradients without artificially distorting the regularization objective.)

\subsection{Sequential Smoothness Regularization ($\mathcal{L}_{\text{smooth}}$)}
\label{subsec:smooth_loss}

While our coordinate sandbox employs a layer-averaging simplification to fit its single-layer expert classifiers, deep sequential architectures (such as foundation models) deploy independent, un-averaged routing coefficients per layer. In these deep sequential networks, allowing layer-wise routing weights to fluctuate wildly and independently across adjacent layers poses a severe representation risk. Because representation spaces are processed sequentially, rapid and incoherent layer-to-layer routing variations introduce a compounding "routing jitter." This jitter results in sequentially compounding representation misalignment that corrupts intermediate activations as they flow through the depth of the network, leading to severe functional degradation.

To theoretically resolve this multi-layer deployment bottleneck, we propose a novel \textbf{Sequential Smoothness Regularizer ($\mathcal{L}_{\text{smooth}}$)}. This constraint penalizes rapid, high-frequency routing weight fluctuations across consecutive layers, encouraging sequential consistency and parameter alignment throughout the depth of the network:
\begin{equation}
\begin{split}
    \mathcal{L}_{\text{smooth}} = \gamma_{\text{smooth}} & \frac{1}{B \cdot (L-1)} \sum_{b=1}^B \sum_{l=2}^L \sum_{k=1}^K \\
    & \left( \alpha_{k, b}(l) - \alpha_{k, b}(l-1) \right)^2
\end{split}
\end{equation}
where $\gamma_{\text{smooth}} \ge 0$ represents a smoothness hyperparameter that controls the regularization strength. Minimizing $\mathcal{L}_{\text{smooth}}$ forces consecutive layers to share similar task ensembling choices, ensuring smooth hidden representation transitions and suppressing sequential routing jitter. Nonetheless, as discussed in Section~\ref{subsec:limitations}, our zero-initialized Softmax prior naturally serves as an inherent, implicit smoothness regularizer, holding layer-by-layer routing weights close to their uniform, maximum-entropy state and preventing routing jitter from emerging, even in the absence of explicit, hyperparameter-sensitive smoothness losses.

\subsection{Multi-Objective Calibration and Vectorized Assembly}
The complete trainable parameters of the router (weights $W \in \mathbb{R}^{L \times K \times d}$ and biases $B \in \mathbb{R}^{L \times K}$ in our framework, including both the well-regularized standard Softmax baseline and its variance-regularized VR-Router variant) are optimized end-to-end on simulated calibration batches. Crucially, we perform \textbf{true sample-specific parameter assembly during training itself}, completely avoiding any batch-averaging compromise or training-time collapse. For a calibration batch $x$ of size $B$, the merged weights are assembled sample-by-sample:
\begin{equation}
    W_{\text{merged}}^{(l)}(b) = W_{\text{base}}^{(l)} + \sum_{k=1}^K \alpha_{k, b}(l) V_k^{(l)}
\end{equation}
The cross-entropy loss $\mathcal{L}_{CE}$ is computed directly on these sample-specific merged weights using highly optimized vectorized operations:
\begin{equation}
    \mathcal{L}_{CE} = -\frac{1}{B} \sum_{b=1}^B \log P(y_b \mid x_b; W_{\text{merged}}(b))
\end{equation}
The complete multi-objective loss is defined as:
\begin{equation}
    \mathcal{L}_{\text{total}} = \mathcal{L}_{CE} + \mathcal{L}_{\text{reg}} + \mathcal{L}_{VR} + \mathcal{L}_{\text{smooth}}
\end{equation}
where $\mathcal{L}_{\text{reg}}$ represents $L_2$ routing weight decay:
\begin{equation}
    \mathcal{L}_{\text{reg}} = \lambda_{wd} \sum_{l=1}^L \sum_{k=1}^K \left( \|W_{l, k}\|_2^2 + B_{l, k}^2 \right)
\end{equation}

\subsection{Implicit Regularization via Zero-Initialized Softmax}
A critical finding of our empirical audit is the profound role of \emph{zero-initialization} in the routing layers. By initializing the weights $W$ and biases $B$ to exact zeros, the router initially outputs equal, uniform coefficients for all task experts ($\alpha_{k, b}(l) = 1/K$ for all $k$). This zero-initialization acts as a maximum-entropy uniform prior that naturally minimizes intra-task variance and restricts the model to a highly-regularized, convex compromise.

Because the calibration split is extremely small ($|D_{\text{cal}}| = 64$) and training is short (100 epochs), the router weights remain constrained close to their initial zero state under $L_2$ decay. This zero-initialized Softmax architecture acts as a powerful \textbf{implicit task-variance regularizer}. It prevents the routing weights from overfitting to high-frequency sample noise, ensuring that predicted sample-wise routing coefficients remain remarkably stable and consistent. This implicit regularization explains the flat sensitivity curves and why our zero-initialized routing models completely avoid the severe \emph{Vectorization Collapse} observed in random-initialized models (which drop to 42.34\% under $B=1$ vectorized deployment).

\subsection{Resolving the Conceptual Premise of Dynamic Routing}
We explicitly address the conceptual role of Task-Variance Regularization ($\mathcal{L}_{VR}$). One might argue that minimizing intra-task variance ($\sigma^2_{\text{task}} \to 0$) forces the router to act as a static task classifier, rendering the dynamic routing machinery redundant. 

However, we contend that $\mathcal{L}_{VR}$ acts as a soft-clustering constraint that filters out high-frequency input noise while preserving coarse task boundaries. Under data scarcity, unregularized dynamic routers suffer from severe overfitting, misinterpreting sample noise as task indicators. Minimizing intra-task variance prevents this overfitting, forcing the router to align its predictions with robust task subspaces. At the same time, the zero-initialized Softmax architecture allows the router to dynamically scale task coefficients smoothly across tasks on boundary samples, achieving a robust compromise that a rigid, discrete task classifier cannot realize.

During deployment, we bypass batch-averaging entirely. We employ PyTorch's vectorized mapping API (`torch.vmap`) or highly optimized einsum operations to perform sample-specific, batch-independent forward passes. This allows each input sample in a mixed batch to be processed by its own dynamically assembled parameters, realizing true batch-independent multi-task routing without any batch-averaging bottlenecks.

\subsection{Systems-Level Complexity, Hardware Bottlenecks, and Scaling to Billion-Parameter Models}
\label{subsec:systems_complexity}
While sample-specific vectorized assembly mathematically resolves Heterogeneity Collapse and Vectorization Collapse by decoupling routing from the batch size, it introduces substantial hardware-level memory and computational challenges. At the physical layer, constructing $B$ distinct, sample-specific weight matrices of size $D_{\text{in}} \times D_{\text{out}}$ scales the active model parameter footprint by a factor of $B$. For deep architectures with hundreds of millions or billions of parameters, storing $B$ separate replicas of the network's weights in GPU memory is extremely memory-intensive, often leading to severe out-of-memory (OOM) errors during both training and inference.

Furthermore, this sample-wise vectorized assembly transforms standard, highly-optimized General Matrix Multiplications (GEMMs) into batch-wise operations (such as \texttt{torch.einsum('b...d,b...o->b...o')}). These operations suffer from a massive drop in arithmetic intensity (the ratio of floating-point operations to memory access). Because the processor must fetch $B$ distinct weight matrices from High Bandwidth Memory (HBM) instead of reusing a single static weight matrix across the entire batch, the execution becomes heavily memory-bandwidth bound rather than compute bound. On modern hardware accelerators, this memory bandwidth bottleneck leads to severe processing stalls and a significant increase in wall-clock latency compared to static Uniform Merging (which has zero runtime overhead).

To mitigate these systems-level limitations when transitioning from our controlled sandbox to multi-billion parameter foundation models (such as LLaMA-70B or CLIP-ViT-L/14), we propose that \emph{low-rank parameter assembly} becomes a mandatory structural requirement rather than an optional optimization. By applying dynamic routing coefficients exclusively to low-rank parameter-efficient adapters (such as LoRA \cite{hu2021lora}), where adapter weights are decomposed into low-rank matrices $A \in \mathbb{R}^{D_{\text{in}} \times r}$ and $B \in \mathbb{R}^{r \times D_{\text{out}}}$ with rank $r \ll \min(D_{\text{in}}, D_{\text{out}})$, the vectorized assembly is restricted to less than $1\%$ of the total network parameters. In this low-rank regime, the active memory footprint is virtually unaffected, and the memory bandwidth overhead of fetching sample-specific adapter weights is small enough to be easily hidden behind the compute-heavy base model forward pass, realizing a highly scalable, systems-efficient deployment pipeline.
